% =============================================================================
%  Liste des acronymes
%
%  Volontairement composée avec un simple `longtable` plutôt qu'avec le paquet
%  `acronym` ou `glossaries` : le modèle INPT ne les charge pas nécessairement,
%  et une liste statique ne peut pas se désynchroniser d'un compteur.
% =============================================================================

\chapter*{Liste des acronymes}
\addcontentsline{toc}{chapter}{Liste des acronymes}
\markboth{LISTE DES ACRONYMES}{LISTE DES ACRONYMES}

\begin{longtable}{@{}p{2.6cm}p{11.4cm}@{}}
\toprule
\textbf{Acronyme} & \textbf{Signification} \\
\midrule
\endfirsthead
\toprule
\textbf{Acronyme} & \textbf{Signification} \\
\midrule
\endhead

API & \emph{Application Programming Interface} — interface de programmation
applicative \\
BAM & Bank Al-Maghrib — banque centrale du Royaume du Maroc \\
BVC & Bourse des Valeurs de Casablanca \\
DCC & \emph{Dynamic Conditional Correlation} — corrélation conditionnelle
dynamique (Engle, 2002) \\
DSR & \emph{Deflated Sharpe Ratio} — ratio de Sharpe déflaté \\
DVC & \emph{Data Version Control} — outil de versionnement des données et des
pipelines \\
EM & \emph{Expectation-Maximization} — algorithme d'estimation employé par le
modèle de Markov caché \\
ETF & \emph{Exchange-Traded Fund} — fonds indiciel coté \\
EWMA & \emph{Exponentially Weighted Moving Average} — moyenne mobile à
pondération exponentielle \\
FRED & \emph{Federal Reserve Economic Data} — base de données macroéconomiques
de la Réserve fédérale de Saint-Louis \\
GARCH & \emph{Generalized Autoregressive Conditional Heteroskedasticity} —
modèle de volatilité variable dans le temps \\
HMM & \emph{Hidden Markov Model} — modèle de Markov caché \\
IC & \emph{Information Coefficient} — coefficient d'information \\
INPT & Institut National des Postes et Télécommunications \\
MAD & Dirham marocain (code ISO~4217) \\
ML & \emph{Machine Learning} — apprentissage automatique \\
MPT & \emph{Modern Portfolio Theory} — théorie moderne du portefeuille
(Markowitz, 1952) \\
OOS & \emph{Out-Of-Sample} — hors échantillon \\
PFA & Projet de Fin d'Année \\
REST & \emph{Representational State Transfer} — style d'architecture des
interfaces web \\
RF & \emph{Random Forest} — forêt aléatoire \\
SLSQP & \emph{Sequential Least SQuares Programming} — algorithme
d'optimisation sous contraintes \\
USD & Dollar américain (code ISO~4217) \\
VIX & \emph{Volatility Index} — indice de volatilité implicite du marché
américain \\
XGB & \emph{eXtreme Gradient Boosting} — bibliothèque de \emph{gradient
boosting} d'arbres \\
\bottomrule
\end{longtable}
